\section{Sweeps, activity, and exchange graphs}
\label{sec:finite}

The chain and crossing mechanism below is the one underlying Dey's bound~\cite{Dey1998}. We first convert intervals containing many participating points into
straight-line graphs with controlled crossings.  This gives the classical
exponent $4/3$ and isolates the packing estimate that must be improved.

We first justify the genericity convention of the introduction. The
additional conditions on coordinates, slopes, and proper segment
intersections can be achieved by an arbitrarily small perturbation.  Indeed, the signs
of the finitely many nonzero oriented areas of triples are unchanged by
a sufficiently small perturbation.  Consequently every pair has exactly
the same points on its two sides, and in particular the halving pairs are
preserved.  The additional conditions exclude finitely many proper
algebraic sets of perturbations.

Write $P=\{(x_i,y_i):1\le i\le N\}$ and fix a rank $k$, with
$0\le k\le N$.  At a slope $s$ that is not a determined slope, let
$M_k(s)$ be the set of the $k$ smallest values of $y_i-sx_i$.
Choose one representative slope in each interval between consecutive
determined slopes, including the two unbounded intervals.  These are
the \emph{states}; adjacent states form a \emph{step}.  A step either
leaves $M_k$ unchanged or replaces one point by one other point.
The latter kind of step is a \emph{switch}.  An interval of states
includes its endpoints, and its step set consists of the intervening
steps.  Its \emph{activity} is the number of distinct points whose
membership in $M_k$ changes at one of those steps.  Thus activity counts
points, rather than switches.  For two sets of size $k$ put
\[
 d(M,M')=|M\setminus M'|=\tfrac12|M\mathbin{\triangle}M'|.
\]
Activity records membership changes throughout the interval, whereas
$d(M,M')$ records only the difference of its endpoint sets.
This is a metric, and one step changes distance from a fixed state by
at most one. At middle rank, the same endpoint-set distance occurs
in Pinchasi's study of slope-ordered balanced cuts and measure
concentration~\cite{Pinchasi2009}; the activity reduction below also
handles points that return to their initial membership.

At a switch the gained point has larger first coordinate than the lost
point: the difference of their two projections changes sign as $s$
increases.  The sum of first coordinates in $M_k$ therefore increases
strictly at every switch.  In particular an interval of positive
activity has different endpoint sets.

Our main intermediate result strengthens the classical crossing bound
for collections of disjoint sweep intervals. Its quantifier is uniform
in both the rank and the activity threshold.

\begin{theorem}[Uniform activity packing]\label{thm:activity-packing}
For every $\xi>0$ there is $Q_\xi<\infty$ such that, for every generic
$N$-point configuration, every rank, and every integer $1\le b\le N$
with $N/b\ge Q_\xi$, any collection of intervals with disjoint step
sets and activity at least $b$ has at most $\xi(N/b)^{4/3}$ members.
\end{theorem}

We first establish the finite reductions and the corresponding
$O((N/b)^{4/3})$ bound. Sections~\ref{sec:rank-limits}--\ref{sec:boundary}
prove the improvement to the displayed arbitrarily small constant.

\begin{lemma}[Activity and endpoint distance]\label{lem:activity-distance}
Fix a generic $N$-point set and a rank $k$.  For an integer $b\ge4$,
any collection of intervals of activity at least $b$ can be divided
into two classes.  At rank $k+\lfloor b/4\rfloor$ the intervals in
the first class have endpoint distance at least $\lfloor b/4\rfloor$;
at rank $k-\lfloor b/4\rfloor$ the intervals in the second class do
so.  Each rank used by a nonempty class is feasible.  For $1\le b<4$,
the original rank already gives endpoint distance at least one.
\end{lemma}

\begin{proof}
Set $c=\lfloor b/4\rfloor$.  Classify an interval according to whether
at least $b/2$ of its active points were unmarked initially; put all
remaining intervals in the second class.  In the first case let $E$
be the initially marked set at rank $k+c$.  At most $c$ of the initially
unmarked active points lie in $E$, so at least
$\lceil b/2\rceil-c\ge c$ of them lie outside $E$.

Take such a point $u$.  At some state it enters rank $k$, and is then
below at least $c+1$ of the $k+c$ points of $E$ in projection order.
Initially it was above all of $E$.  For each of those $c+1$ points the
projection difference is affine in the slope, so the new comparison
persists to the final state.  Let $E'$ be the final rank-$(k+c)$ set.
If at least $c$ of the specified points outside $E$ lie in $E'$, then
$|E'\setminus E|\ge c$.  Otherwise one of them, say $u$, is outside
$E'$.  All $c+1$ points of $E$ that remain above $u$ are then outside
$E'$, so again the endpoint distance is at least $c$.
There were at least $\lceil b/2\rceil$ initially unmarked points, which
also proves $k+c\le N$.

The other class has at least $b/2$ initially marked active points.
Apply the preceding argument to the complementary marks and the
negated projections; enlarging the complementary rank by $c$ is
the same as reducing the original rank by $c$.  Its feasibility follows
in the same way.  The case $b<4$ follows from the strict increase of
the sum of first coordinates at switches.
\end{proof}

\begin{lemma}[Greedy batches]\label{lem:greedy-batches}
Fix a rank and an integer $d\ge1$, and put $a=\lceil d/2\rceil$.
Starting at the first state, end a batch at the first subsequent state
at distance $a$ from its start, and repeat while possible.  If $R$
complete batches are obtained, every collection of intervals with
disjoint step sets and endpoint distance at least $d$ has at most
$R$ members.
\end{lemma}

\begin{proof}
Distances change by at most one at a step, so every completed batch has
distance exactly $a$.  If an interval contains no batch end strictly
after its start, both its endpoints are at distance at most $a-1$
from the same preceding batch start.  Their mutual distance is at most
$2(a-1)<d$, a contradiction.  This applies also in the final unfinished
batch.  Assign to each interval a batch end strictly after its start
and at or before its end.  Disjoint step sets cannot receive the same
batch end.
\end{proof}

Now consider any consecutive batches of endpoint distance $a$.  In a
batch let $L$ be the $a$ points lost between its endpoint states and
$G$ the $a$ points gained.  Draw all $a^2$ segments from $L$ to $G$.
Every segment points strictly to the right, and its slope lies
strictly between the two state slopes of the batch.  Indeed, for
$u\in L$ and $v\in G$, their projection order is reversed between
the two states.  A pair cannot occur in two batches because its
determined slope would have to belong to two disjoint open slope
intervals.  The resulting exchange graph is therefore simple.

\begin{lemma}[Two chain partitions]\label{lem:chain-partitions}
For a consecutive window of batches, let $S$ be the number of distinct
endpoints of its exchange edges.  Those edges partition into at most
$aS$ convex polygonal chains, and also into at most $aS$ concave
polygonal chains.  The number of properly crossing edge pairs is at
most $2a^2S^2$.
\end{lemma}

\begin{proof}
Decompose each complete bipartite exchange into $a$ perfect matchings,
and number the matchings from $1$ to $a$.  For one matching number,
place a token at each marked point at the start of the window.
In every batch move the token at each lost point to its matched gained
point, leaving all other tokens in place.  Along a token's moving
segments the first coordinate increases, and segment slopes increase
strictly from batch to batch.  Its nonempty trajectory is a convex
polygonal chain.  A token that ever moves starts at an endpoint in
the window: its first move otherwise could not occur.  Thus there
are at most $S$ nonempty chains for each matching number.

Instead place tokens at the unmarked points and move them from gained
to lost points along the same matchings.  They move left as time
increases.  Reading a trajectory from left to right reverses the
batch order, so its slopes decrease.  This gives the concave
partition, again with at most $aS$ chains.

On their common horizontal domain the difference of a convex chain
and a concave chain is a convex piecewise affine function.  It has
at most two sign-changing zeros, and hence the chains have at most
two proper transverse intersections in edge interiors.  Shared
segments do not count as proper crossings and do not increase this
bound.  Assign one edge of each crossing pair to its convex chain
and the other to its concave chain.  The genericity assumptions
exclude multiple proper crossings at one point.  Summing over chain
pairs proves the stated bound.
\end{proof}

The classical crossing lemma is due independently to Ajtai, Chv{\'a}tal,
Newborn and Szemer{\'e}di, and to Leighton
\cite{AjtaiEtAl1982,Leighton1983}. The asymmetric sampling argument
below is the imbalanced bipartite crossing argument of Pach, Solymosi
and Tardos~\cite[Lemmas~2.1--2.2]{PachSolymosiTardos2010}; we include
the reduction to designated endpoints and the constants needed here.
We need a version of the crossing estimate that also controls edges
whose endpoint lies in a prescribed small set.  In the next lemma
$\operatorname{cr}(G)$ denotes the number of properly crossing edge
pairs in the given drawing, not the minimum over drawings.

\begin{lemma}[Crossings with designated endpoints]
\label{lem:designated-crossing}
Let a simple graph be drawn with straight segments on $v$ distinct
points, with no vertex in an edge interior and no crossings of
incident edges.  Suppose every edge has an endpoint in a specified
set of at most $m\le v$ vertices.  If the graph has $e$ edges, then
\[
 e^3\le 2304\,mv\,\operatorname{cr}(G)+13824\,v^3.
\]
In particular, if $e\ge24v$, then
$e^3\le2304\,mv\,\operatorname{cr}(G)$.
\end{lemma}

\begin{proof}
The assertion is immediate if $e<24v$ or $e=0$.  Select each
designated vertex independently with probability $1/2$, put these
selected vertices on one side, and put all other vertices on the
other side.  Every edge crosses this cut with probability $1/2$.
Some resulting bipartite subgraph therefore has $e_0\ge e/2$ edges,
with side sizes at most $m$ and $v$.

Retain vertices on its two sides independently with probabilities
\[
 p=\frac{12v}{e_0},\qquad q=\frac{12m}{e_0},
\]
respectively.  Both probabilities lie in $[0,1]$.  Any graph in this
drawing satisfies $e'\le3v'+\operatorname{cr}'$: deleting at most
one edge for each crossing leaves a simple plane graph.  Taking
expectations gives
\[
 pq e_0\le3(pm+qv)+p^2q^2\operatorname{cr}(G).
\]
The factor $p^2q^2$ is valid because crossing edges have four
distinct endpoints.  Substitution yields
$e_0^3\le288mv\operatorname{cr}(G)$, and $e_0\ge e/2$ proves
the dense assertion.  The additional term $24^3v^3$ covers the
remaining case.
\end{proof}

For $R$ consecutive batches on $N$ points the graph has $Ra^2$ edges
and at most $2a^2N^2$ crossing pairs.  Applying
Lemma~\ref{lem:designated-crossing} with $m=v\le N$, and separating
the sparse case, gives
\[
 R\le C\left(\frac{N}{a^2}+\left(\frac Na\right)^{4/3}\right)
 \le C'\left(\frac Na\right)^{4/3}.
\]
Here $a\ge1$ and $N\ge1$, so the first summand is absorbed by the
second.  Lemmas~\ref{lem:activity-distance} and
\ref{lem:greedy-batches} now show that the number of intervals with
disjoint step sets and activity at least $b$ is
$O((N/b)^{4/3})$, uniformly in the rank.  The reductions also show
that a uniform improvement to $o((N/a)^{4/3})$ for consecutive
distance-$a$ batches, as $N/a\to\infty$, gives the same little-$o$
improvement for activity packing.  There are only two rank classes,
and for $b\ge4$ their batch parameter obeys $b/16\le a\le b/4$;
for $b<4$ it is $a=1$.

We record also the localized bound to be used below.  If a window
has $S$ distinct endpoints and $e_F$ of its edges have their left
endpoint among a specified $m$ of those endpoints, then
\begin{equation}\label{eq:finite-localized-crossing}
 e_F^3\le4608\,a^2mS^3+13824\,S^3.
\end{equation}
Indeed this subgraph has at most $S$ vertices and inherits the
crossing bound from Lemma~\ref{lem:chain-partitions}.  The second
term explicitly retains the sparse case; it may only be discarded
after an appropriate limiting normalization.



The remainder of Sections~\ref{sec:rank-limits}--\ref{sec:boundary}
proves Theorem~\ref{thm:activity-packing}. The finite reductions show that a
failure produces a sequence of consecutive distance-$a$ batches with
$N/a\to\infty$ and $R\asymp(N/a)^{4/3}$. We will derive a contradiction
from such a sequence. Afterward Section~\ref{sec:recurrence} converts
this uniform statement into the power saving of
Theorem~\ref{thm:halving-power}.
